\section{Future work: people true to the person}\label{future-work-people-true-to-the-person}

The minimal sim body is a start, and what follows are goals rather than results.

\textbf{Bodies true to the person.} The sim body is to become fully musculoskeletal and true to the person, older people and a curved spine included, instead of a resized template. Joint centres are to come from bone landmarks on its SOMA mesh, then mass, muscles, joint ranges, spine and bone shape, each input SOMA lacks marked as assumed with a range. A full musculoskeletal body with 416 muscles already runs on Newton but still drifts 1.5 cm from plain MuJoCo after 300 steps, so Newton is to be validated for bodies first. Its joint angles, then its joint moments and ground forces, are to be held to OpenCap's LabValidation data.

\textbf{Scenes from motion.} Motion comes from the SOMA ecosystem, read from video by GEM-X \citep{gemx} or made from sentences by Kimodo \citep{rempe2026kimodo}, without limit. The ground a motion implies becomes a SCORE scene, with no video and nothing that identifies a person. A sim body follows the motion on that ground, and physics computes the ground contact, the joint torques and the forces of its 416 muscles. The force that neither the ground nor the muscles can supply is the residual. It is recorded and never applied, and only motion with a small residual is kept. Biomechanics uses the same residual to judge recorded motion \citep{delp2007opensim, werling2023addbiomechanics}, and PhysibleMDM \citep{stolzle2026physiblemdm} lowers it for generated robot motion; we use it on generated motion in a muscle-driven body.

\textbf{Layered training sets.} Each run also yields layered training sets, in which a rendered picture of a clothed person in a scene is paired with the exact body, joint angles, ground forces, muscle activations and forces, and residual beneath it. From these a fast model can be trained to detect muscle metrics from video, and specialised, for example for clinical use, because its training data comes from the simulation. BioHuman10M \citep{huo2026biohuman} pairs rendered video with 112-muscle activations and estimated ground forces; our sets add scenes made for the purpose, ground forces from simulated contact, a 416-muscle body and a residual per sample. Whether such a model carries over to real force-plate and muscle recordings remains to be shown.

\textbf{Brains.} A brain is a trained model that drives a sim body's muscles. Trained on low-residual motion and placed in new worlds, it adds movement no recording holds. KINESIS \citep{simos2025kinesis} and MuscleMimic \citep{li2026musclemimic} train muscle-driven controllers on motion capture; to our knowledge none has been trained on generated motion filtered by its residual, a claim that still needs a comparison with both.

\textbf{The data loop.} The more compute behind the loop, the more data it makes, which is why SCORE is open, so others can add their compute and share the sets back. Models for clinical use are trained outside SCORE, on data that never enters it.
